\clearpage
\hypertarget{native-481-january4}{}
\begin{center}
\rule{0.55\linewidth}{0.4pt}\\[2pt]
Astronomische Nachrichten. Band XXVI. Nr. 623. S. 367. 17. Januar 1848.\\[2pt]
\rule{0.55\linewidth}{0.4pt}
\end{center}
\begin{center}[\textsc{Gauss} an \textsc{Schumacher}.]\quad Göttingen, 4. Januar 1848.\end{center}

Nach langer Unterbrechung klarte der Himmel gestern auf, wenn auch nur kurz. So konnte ich wieder eine Meridianbeobachtung der \emph{Flora} machen, die ich Ihnen hier mitteile.

\begin{center}
Beobachtung der Flora.\\[3pt]
\setlength{\tabcolsep}{4pt}
\begin{tabular}{l|r|r|r}
1848 & Mittl. Zeit in Göttingen & Rektaszension & Deklination \\\hline
Januar 3 & $9^{\rm u}\ 20^{\rm m}\ 4^{\rm s}2$ & $62^\circ\ 37'\ 13''8$ & $+15^\circ\ 59'\ 24''6$ \\
\end{tabular}
\end{center}

Die Elemente von Professor \textsc{Goldschmidt} weichen davon ab (nach seiner Rechnung) um

\begin{center}
$+18''2$ in der Rektaszension.\qquad $+2''4$ in der Deklination.
\end{center}

